% \documentclass[../Main.tex]{subfiles}

% \begin{document}
\label{chapter3:Our Threshold Batch Decryption system}

\section{Technical Overview}
At a high level, our scheme adapts the toolbox for thresholdizing lattice gadget
trapdoors from~\cite{AC:ALLW25} to the batched IBE (BIBE) construction of
Boneh, Laufer and Tas~\cite{EPRINT:BonLauTas25}. In~\cite{EPRINT:BonLauTas25},
the authors use the shifted multi-preimage trapdoor sampler of Waters, Wee, and
Wu~\cite{EC:WatWeeWu25} as an accumulator that binds all identities within a
batch. As a result, the pre-decryption key is a Gaussian preimage of a masking
vector, in the same spirit as the dual-Regev scheme. By applying the partial
trapdoor toolbox, we let multiple parties independently sample fragments of this
preimage and recombine them afterwards, thereby removing the single point of
failure in the original decryption authority.

Throughout, $N$ denotes the number of parties (the committee) and $k\le N$ the
reconstruction threshold, so that any $k$-out-of-$N$ quorum can decrypt while any
$k-1$ parties learn nothing.

\paragraph{3.1.1. The Index-Dependent Batch IBE scheme by Boneh, Laufer and Tas~\cite{EPRINT:BonLauTas25}.}
Let $r_1, \dots, r_\ell$ be the batch of identity tags, where $r_i$ corresponds to
index $i$. The pre-decryption algorithm first samples a preimage vector
$
  \mathbf{v} = \mathsf{col}(\mathbf{v}^1_1, \dots, \mathbf{v}^1_\ell,
  \hat{\mathbf{c}}, \mathbf{v}^2_1, \dots, \mathbf{v}^2_\ell)
$
of $(\mathbf{u}, \dots, \mathbf{u}) \in \mathbb{Z}_q^{n\ell}$ under
$[\mathbf{A} \mid \hat{\mathbf{A}}]$, so that
\[
  [\mathbf{A}_i \mid \mathbf{G} \mid H_{\mathsf{ro}}(r_i)]
  \begin{bmatrix} \mathbf{v}^1_i \\ \hat{\mathbf{c}} \\ \mathbf{v}^2_i \end{bmatrix}
  = \mathbf{u}
  \quad \text{for all } i = 1, \dots, \ell.
\]
Given only $\mathsf{pk}$, this preimage is sampled publicly by generating a
trapdoor $\mathbf{T}_{\mathbf{A}}$ for $\mathbf{A}$ and calling
$\mathsf{SampleLeft}(\mathbf{A}, \hat{\mathbf{A}}, \mathbf{T}_{\mathbf{A}},
(\mathbf{u}, \dots, \mathbf{u}), \sigma)$. To enable decryption of all $\ell$
ciphertexts, the algorithm then returns the pre-decryption key
\[
  \mathsf{sbk} \leftarrow \mathsf{SamplePre}(\mathbf{C}, \mathbf{T}_{\mathbf{C}},
  \mathbf{c} := \mathbf{G}\hat{\mathbf{c}}, \sigma),
  \qquad \text{so that } \mathbf{C} \cdot \mathsf{sbk} = \mathbf{G}\hat{\mathbf{c}}.
\]
The triple $(\mathsf{sbk}, \mathbf{v}^1_i, \mathbf{v}^2_i)$ is then an IBE secret
key for the tag $r_i$, satisfying
\[
  [\mathbf{C} \mid \mathbf{A}_i \mid H_{\mathsf{ro}}(r_i)]
  \begin{bmatrix} \mathsf{sbk} \\ \mathbf{v}^1_i \\ \mathbf{v}^2_i \end{bmatrix}
  = \mathbf{u}.
\]
For decryption to succeed, every decryptor must obtain the \emph{same} preimage
$\mathbf{v}$; this is ensured by deriving the randomness of $\mathsf{SampleLeft}$
from a hash function modeled as a random oracle.

\paragraph{3.1.2. Lattice trapdoor secret sharing technique in~\cite{AC:ALLW25}.}
The key idea is to secret-share in the \emph{image space} rather than the
preimage space, which sidesteps the norm-growth problems of secret-sharing short
vectors. Recall that the gadget matrix $\mathbf{G} = \mathbf{I}_n \otimes
\mathbf{g}^\top$, with $\mathbf{g}^\top = (1, 2, \dots, 2^{\tilde{q}-1})$, admits a
publicly known short basis and hence allows preimage sampling without any secret;
a \emph{gadget trapdoor} (or $\mathbf{G}$-trapdoor) of a matrix $\mathbf{A}$ is then
a short matrix $\mathbf{U}$ with $\mathbf{A}\mathbf{U} \equiv \mathbf{G} \bmod q$
~\cite{C:MicPei12}, which reduces preimage sampling for $\mathbf{A}$ to the easy
gadget case. Now let $\mathbf{V} \in \mathbb{Z}_q^{k\times N}$ be the Vandermonde
matrix whose $j$-th column is
$\mathbf{v}_j = (1, v_j, \dots, v_j^{k-1})^\top$, corresponding to $k$-out-of-$N$
Shamir sharing over the evaluation points $(v_j)_{j\in[N]}$. Each party's partial
trapdoor is a short matrix $\mathbf{U}_j \in \mathbb{Z}_q^{mk\times m}$ satisfying
\[
  \mathbf{A}\cdot\mathbf{U}_j \equiv \mathbf{v}_j \otimes \mathbf{G} \bmod q,
  \qquad \mathbf{A} \in \mathbb{Z}_q^{nk\times mk},
\]
i.e.\ a gadget trapdoor of $\mathbf{A}$ tagged by $\mathbf{v}_j$, so that
$\mathbf{U}_j$ can only sample preimages whose image lies in the
$\mathbb{Z}_q$-subspace spanned by $\mathbf{v}_j \otimes \mathbf{I}_n$; in effect,
the image space is partitioned into $N$ subspaces, one per party. To sample a
preimage of a target $\mathbf{y}$ for a recovery set $T \subseteq_k [N]$, one
decomposes $\mathbf{y} \equiv \sum_{j\in T} \mathbf{v}_j \otimes \mathbf{z}_j
\bmod q$, which is possible since any $k$ columns of $\mathbf{V}$ are invertible.
Party $j$ then locally produces a partial preimage $\mathbf{x}_j$ with
$\mathbf{A}\mathbf{x}_j \equiv \mathbf{v}_j \otimes \mathbf{z}_j \bmod q$, and the
full preimage is simply the sum $\mathbf{x} = \sum_{j\in T} \mathbf{x}_j$, for which
$\mathbf{A}\mathbf{x} \equiv \mathbf{y} \bmod q$. Moreover, rather than outputting
the insecure $\mathbf{x}_j := \mathbf{U}_j \cdot \mathbf{G}^{-1}(\mathbf{z}_j)$
(which leaks $\mathbf{U}_j$ by linear algebra after enough queries), one samples a
Gaussian $\mathbf{d}_j$ with $\mathbf{G}\mathbf{d}_j \equiv \mathbf{z}_j \bmod q$
and covariance shaped by $\mathbf{U}_j$, so that $\mathbf{x}_j = \mathbf{U}_j
\mathbf{d}_j$ is a spherical Gaussian over the rank-$m$ sublattice
$\Lambda(\mathbf{U}_j)$, leaking only $\Lambda(\mathbf{U}_j)$ and not $\mathbf{U}_j$
itself.

\paragraph{3.1.3. Our threshold batch identity-based encryption scheme.}
The setup algorithm now samples, for each party $j\in[N]$, a partial trapdoor tuple
$\mathsf{ptd}_j = (\mathbf{U}_j, \mathbf{C}_j, \mathbf{T}_{\mathbf{C}_j})$ via
$\mathsf{PTrapGen}(\mathbf{C}, \mathbf{T}_{\mathbf{C}}, j)$, which are then secretly
distributed to the corresponding parties. With these in hand, the pre-decryption
algorithm proceeds exactly as in~\cite{EPRINT:BonLauTas25} up to the point of
forming the shifted vector $\mathbf{c} := \mathbf{G}\hat{\mathbf{c}}$. The crucial
difference is that the single $\mathsf{SamplePre}$ call producing $\mathsf{sbk}$ is
replaced by its distributed analogue: each party $j$ in a quorum
$T = \{v_1,\dots,v_k\} \subseteq_k [N]$, where every $v_i\in\mathbb{Z}$ is the
identification value pre-assigned to a distributed party, locally samples a partial
pre-decryption key
\[
  \mathbf{sbk}_j \leftarrow \mathsf{PSampPre}(\mathbf{C}, \mathsf{ptd}_j, T,
  \mathbf{c} := \mathbf{G}\hat{\mathbf{c}}, \sigma_2) \in \mathbb{Z}_q^{2mk}.
\]
Once partial keys from a full quorum are available, anyone can reconstruct the
complete pre-decryption key by simple summation,
\[
  \mathbf{sbk} \leftarrow \mathsf{Rec}\big((\mathbf{sbk}_j)_{j \in T}\big)
  \in \mathbb{Z}_q^{2mk},
\]
by the correctness of the partial trapdoor scheme, is satisfied
$\mathbf{C}\cdot\mathbf{sbk} = \mathbf{G}\hat{\mathbf{c}}$ with $\|\mathbf{sbk}\|$
bounded. The remainder of the procedure (encryption and the final rounding-based
decryption) is identical to the BIBE scheme of~\cite{EPRINT:BonLauTas25}. In this
way the decryption authority is split among $N$ parties, and no coalition of fewer
than $k$ of them can recover a pre-decryption key.

\paragraph{3.1.4. Security.}
The security argument largely mirrors that of the BIBE construction
in~\cite{EPRINT:BonLauTas25}, with a few subtle differences introduced by the
trapdoor secret-sharing layer. First, we work in a selective
TBIBE security game: before observing $\mathsf{pk}$, the adversary must commit to a
corruption set $\mathcal{C}\subseteq_{<k}[N]$ of at most $k-1$ parties, the full
list of pre-decryption (secret-key) queries
$\big(T_i,(r_{i,j})_{j\in[\ell]}\big)_{i\in[L]}$, and the challenge tuple
$(T^\ast, z^\ast, r^\ast)$ specifying the target quorum, the target party, and the
target identity. During the query phase, for the challenge quorum $T^\ast$ the
adversary receives all partial pre-decryption keys \emph{except} the one held by
the target party $z^\ast$, which is withheld. This downgrade from adaptive to selective security is inherited from the partial trapdoor secret-sharing
scheme~\cite{AC:ALLW25}: although its construction is simple, simulating partial
trapdoors and partial preimages without the full trapdoor is delicate and forces
the adversary to declare its corruptions and challenge in advance.

The proof proceeds through a sequence of hybrid games. The early hybrids
reprogram the sampler CRS, the random oracle $H_{\mathsf{ro}}$, and the
explainable-sampler oracle $H_{\mathsf{sp}}$ so that the challenge matrices are
embedded as $\mathbf{A}_{r^\ast} = \mathbf{C}\mathbf{H}_1$ and
$H_{\mathsf{ro}}(r^\ast) = \mathbf{C}\mathbf{H}_2$, while pre-decryption answers
are simulated via $\mathsf{SampleLeft}$ together with the explainability of the
sampler; each step changes the adversary's view by at most a negligible amount.
In the final reduction we build an adversary $\mathcal{B}$ against the
$(nk, 2mk, \chi_0, \chi_1, \sigma_2)$-indistinguishability of the partial
trapdoor, which in turn is implied by the $\kappa$-LWE assumption~\cite{AC:ALLW25}.
The reduction $\mathcal{B}$ forwards the corruption set $\mathcal{C}$ to the
partial-trapdoor challenger, receives the public objects
$(\mathsf{crs}_{\mathsf{samp}}, \tilde{\mathbf{A}}, \tilde{\mathbf{B}},
\mathbf{C}, \mathbf{u})$ together with the corrupt partial trapdoors
$(\mathbf{U}_j)_{j\in\mathcal{C}}$, and on the challenge $(T^\ast, z^\ast)$ obtains
the partial preimages $(\mathbf{x}^\ast_j)_{j\in T^\ast\setminus\{z^\ast\}}$, a
target $\mathbf{y}^\ast$, and challenge components $(\mathbf{ch}_0, ch_1)$. It then
assembles the challenge ciphertext as
\[
  \mathsf{ct}[1] := ch_1 + m_b\cdot\lfloor q/2\rfloor, \qquad
  \mathsf{ct}[2] := \mathbf{ch}_0, \qquad
  \mathsf{ct}[3] := \mathbf{ch}_0^\top\mathbf{H} + \mathbf{e}',
\]
with $\mathbf{e}'\leftarrow\mathcal{D}_{\mathbb{Z},\chi'}^{t+\overline{m}}$. When
$ch_1 = \mathbf{u}^\top\mathbf{s} + e_1$ is a genuine LWE sample, $\mathcal{B}$
perfectly reproduces the hybrid in which the challenge ciphertext is real;
when $ch_1$ is uniform, it reproduces the hybrid in which the ciphertext is
uniformly random. Hence any non-negligible distinguishing advantage of
$\mathcal{A}$ translates directly into an advantage against the partial-trapdoor
indistinguishability, and therefore against $\kappa$-LWE. Combining all the
hybrid transitions, the scheme is selectively secure under the LWE/SIS-type
assumptions underlying the partial trapdoor, with the statistical losses
controlled by Gaussian smoothing and noise-flooding arguments.

\paragraph{3.1.5. The $\kappa$-LWE Hardness Assumption.}
The security of our scheme ultimately rests on a generalized variant of the
Learning With Errors (LWE) problem known as $\kappa$-LWE. We briefly recall the
relevant assumptions here; formal definitions appear in Chapter~2.

\smallskip
\noindent\textit{Learning With Errors.}
Introduced by Regev~\cite{J:Regev09}, who established a quantum reduction from
worst-case lattice problems to its average-case hardness, LWE has since become
the central hardness assumption underlying modern lattice-based cryptography.
The decisional $\mathrm{LWE}_{\mathbb{Z}_q,n,\chi}$ problem is parameterized
by a modulus $q$, a secret dimension $n$, and an error distribution $\chi$ over
$\mathbb{Z}_q$ (typically a discrete Gaussian $\mathcal{D}_{\mathbb{Z},\sigma}$
with small width $\sigma \ll q$). An adversary is given access to an oracle
that either returns noisy linear equations
\[
  (\mathbf{a}_i,\, b_i := \mathbf{a}_i^\top\mathbf{s} + e_i) \in \mathbb{Z}_q^n \times \mathbb{Z}_q,
  \qquad \mathbf{a}_i \leftarrow \mathbb{Z}_q^n,\; e_i \leftarrow \chi,
\]
for a fixed secret $\mathbf{s} \leftarrow \mathbb{Z}_q^n$, or uniformly random
pairs from $\mathbb{Z}_q^n \times \mathbb{Z}_q$. The problem asks to distinguish
between the two cases. In matrix form, given $m$ samples, this is equivalent to
distinguishing $(\mathbf{A},\, \mathbf{A}^\top\mathbf{s} + \mathbf{e})$ from
$(\mathbf{A},\, \mathbf{u})$ for $\mathbf{A} \leftarrow \mathbb{Z}_q^{n\times m}$,
$\mathbf{e} \leftarrow \chi^m$, and $\mathbf{u} \leftarrow \mathbb{Z}_q^m$.

\smallskip
\noindent\textit{From $k$-LWE (LPSS14) to $\kappa$-LWE (ALLW25).}
The $k$-LWE problem was introduced by Ling, Phan, Stehle, and
Steinfeld~\cite{C:LPSS14} in the context of traitor tracing. It augments the
standard LWE oracle with $k$ additional \emph{hint} evaluations of the secret:
for known matrices $\mathbf{R}_i \in \mathbb{Z}_q^{n \times m}$ and fresh noise
$\mathbf{f}_i \leftarrow \chi^m$, the adversary additionally receives
$\mathbf{R}_i^\top\mathbf{s} + \mathbf{f}_i$ for each $i \in [k]$. These hints
make the problem strictly easier for an adversary (more information about
$\mathbf{s}$ is revealed), yet Ling et al.\ proved that $k$-LWE remains hard
under the standard LWE assumption: an adversary breaking $k$-LWE with advantage
$\varepsilon$ yields an adversary breaking standard LWE with advantage
$\Omega(\varepsilon^3/\mathrm{poly}(n))$, where the cubic degradation is
inherent to the self-correcting reduction technique of~\cite{C:LPSS14}.

Albrecht, Lai, Lapiha, and Woo~\cite{AC:ALLW25} generalize this to a
structured $\kappa$-LWE assumption suited to the partial trapdoor setting. Here
the hints are not arbitrary noisy linear functions but rather columns of the
partial trapdoor matrices $\mathbf{U}_i \in \Lambda_q^\perp(\mathbf{A})$
sampled from a structured Gaussian
$\mathcal{D}_{\Lambda_q^\perp(\mathbf{A}),\mathbf{S}_i}$ for prescribed width
matrices $\{\mathbf{S}_i\}_i$. The assumption asserts that
$(\mathbf{A},\,(\mathbf{U}_i)_{i\in[\kappa]},\,\mathbf{A}^\top\mathbf{s}+\mathbf{e})$
is computationally indistinguishable from the same tuple with the last component
replaced by a sample from the $\kappa$-uniform distribution (see Chapter~2).
ALLW25 prove this reduces to standard LWE following the LPSS14 self-correcting
framework, inheriting the same polynomial security loss. In our construction,
the $\kappa = 2m(k-1)$ hint columns visible to an adversary corrupting up to
$k-1$ parties correspond exactly to the partial trapdoor columns satisfying
$\mathbf{C}\mathbf{U}_j \equiv \mathbf{v}_j \otimes \mathbf{G} \pmod{q}$,
placing our scheme directly within the ALLW25 $\kappa$-LWE framework.

\section{Syntax}


A threshold batch identity-based encryption (\(\mathsf{TBIBE}\)) scheme instantiated with a committee of \(N \in \mathbb{N}\) parties, a threshold \(k \in [N]\), a maximum batch size \(\ell \in \mathbb{N}\), a message space \(\mathcal{M}\), subset $T \subseteq_k [N]$ and an identity space \(\mathcal{I}\) consists of five PPT algorithms:
\begin{itemize}
    \item
    $
    \mathsf{TBIBE.Setup}(1^\lambda,1^\ell,1^k, 1^N)\rightarrow (\mathsf{pk},\mathsf{sk}_1,\ldots,\mathsf{sk}_N).
    $

    \item
    $
    \mathsf{TBIBE.Enc}(\mathsf{pk},r,m)\rightarrow \mathsf{ct}.
    $

    \item
    $
    \mathsf{TBIBE.PreDec}(\mathsf{sk}_i,T, (r_i)_{i\in[\ell]})\rightarrow \mathsf{sbk}_i.
    $

    \item
    $
    \mathsf{TBIBE.Combine}(\mathsf{pk},T\subseteq [N],(\mathsf{sbk}_i)_{i\in T})\rightarrow \mathsf{sbk}.
    $

    \item
    $
    \mathsf{TBIBE.Dec}(\mathsf{pk},\mathsf{sbk},(\mathsf{ct}_i)_{i\in[\ell]})\rightarrow (m_i)_{i\in[\ell]}.
    $
\end{itemize}

\begin{definition}[Correctness]
\label{def:TBIBE_correctness}
A $(k,N)$-TBIBE is said to be correct if, for any identity set
$(r_i)_{i\in [\ell]} \in \mathcal{I}^\ell$, and $(m)_{i \in [\ell]} \in \mathcal{M}^\ell$, subset $T \subseteq_t [k]$
it holds that
\[
\Pr\left[
\begin{array}{c|l}
(m_i)_{i \in [\ell]} = (m_i')_{i \in [\ell]}
&
\begin{array}{l}
\mathbf{sbk}_j \leftarrow \mathsf{TBIBE.PreDec}( \mathsf{sk}_j, T, (r_i)_{i\in [\ell]}),
\quad \forall j \in T \\[2mm]
\mathsf{c}_i \leftarrow \mathsf{TBIBE.Enc}(\mathsf{pk}, r_i, m_i)\\[2mm] 
\mathsf{sbk} \leftarrow \mathsf{TBIBE.Combine}(\mathsf{pk},T\subseteq [N],(\mathsf{sbk}_i)_{i\in T})\\[2mm]
m'_i \leftarrow \mathsf{TBIBE.Dec}\bigl(\mathsf{pk},(\mathsf{sk}_j)_{j \in T}, (\mathsf{c}_i)_{i \in \ell}\bigr)
\end{array}
\end{array}
\right]
\geq 1 - \mathsf{negl}(\lambda).
\]
\end{definition}

\begin{definition}[Security]
\label{def:TBIBE_security}

A $(k,N)$-TBIBE scheme
$
\Pi_{\mathsf{TBIBE}} = (\mathsf{Setup}, \mathsf{Enc}, \mathsf{PreDec}, \mathsf{Combine}, \mathsf{Dec})
$
is said to be selectively secure, if for any PPT, $\mathcal{A}$
it holds that
\[
\left|
\Pr\left[\mathrm{Exp}^{0}_{\Pi_{\mathsf{TBIBE}},\mathcal{A}}(1^\lambda)=1\right]
-
\Pr\left[\mathrm{Exp}^{1}_{\Pi_{\mathsf{TBIBE}},\mathcal{A}}(1^\lambda)=1\right]
\right|
< \mathrm{negl}(\lambda),
\]
where $\mathrm{Exp}^{b}_{\Pi_{\mathsf{TBIBE}},\mathcal{A}}$ is defined in Fig.~\ref{fig:TBIBE-security-experiment} and $\kappa$ is negligible.
\end{definition}

\section{Construction}
We introduce our $k\text{-out-of-}N$ threshold batch decryption scheme. The construction is parameterized followed Table~\ref{tab:parameters}

\begin{definition}[The hash function $H$]
\label{def:hash}
The hash function
$H : \mathcal{I} \to \{0,1\}^{d}$
is an injective function that maps an identity tag
$r \in \mathcal{I}$
to a bit vector in
$\{0,1\}^{d}$


\end{definition}

\begin{definition}[The hash function $H_{\mathrm{ro}}$ 
\label{def:hashro}
\cite{EPRINT:BonLauTas25}]
The hash function
$
    H_{\mathrm{ro}} : \mathcal{I} \to \mathbb{Z}_q^{nk \times \bar{m}}
$
is a random oracle that maps an identity tag $\tau \in \mathcal{I}$ to a matrix in
$\mathbb{Z}_q^{nk \times \bar{m}}$.
\end{definition}

\begin{definition}[The hash function $H_{\mathrm{sp}}$ \cite{EPRINT:BonLauTas25}]
\label{def:hashsp}

Let
$
    \rho_{\mathsf{SL}}
    :=
    \rho(\lambda,\ell nk,\ell t+\widetilde{m},\ell \overline{m},q)
$
denote the randomness length of the explainable $\mathsf{SampleLeft}$ sampler
from Theorem~1 when applied to matrices
$
    \mathbf{A}\in \mathbb{Z}_q^{\ell nk\times(\ell t+\widetilde{m})}
    \text{ and }
    \widehat{\mathbf{A}}\in \mathbb{Z}_q^{\ell nk \times \ell \widetilde{m}} .
$
The hash function
$
    H_{\mathrm{sp}}:
    \mathbb{Z}_q^{\ell nk \times(\ell t+\widetilde{m})}
    \times
    \mathbb{Z}_q^{(\ell t+\widetilde{m})\times \ell \widetilde{m}}
    \times
    \mathcal{I}^{\ell}
    \times
    \mathbb{Z}_q^{nk}
    \times
    \mathbb{Z}^{+}
    \to
    \{0,1\}^{\rho_{\mathsf{SL}}}
$
takes as input a matrix
$\mathbf{A}\in \mathbb{Z}_q^{\ell nk\times(\ell t+\widetilde{m})}$,
a gadget trapdoor
$\mathbf{T}_{\mathbf A}\in
\mathbb{Z}^{(\ell t+\widetilde{m})\times \ell \widetilde{m}}$
for $\mathbf{A}$, identity tags
$(r_1,\ldots,r_\ell)\in\mathcal{I}^{\ell}$,
a vector $\mathbf{u}\in\mathbb{Z}_q^n$, and a Gaussian width parameter
$\sigma$. It outputs the random coins used by $\mathsf{SampleLeft}$.
\end{definition}

\noindent We now describe out threshold decryptoon scheme\\
\vspace{0.4em} 
\noindent$\mathbf{\mathsf{TBIBE.Setup}}(1^\lambda,1^\ell,1^k, 1^N):$ The trusted server calculates:
\[
\left\{
\begin{aligned}
&\mathsf{crs}_{\mathsf{samp}} \leftarrow \mathsf{Gen}(1^\lambda,1^d),
\mathbf u \leftarrow\mathbb Z_q^{nk},\\
&(\mathbf C,\mathbf{T}_{\mathbf{C}}) \leftarrow \mathsf{TrapGen}(1^{nk},q,2mk)\\
\end{aligned}
\right.
\]

For each party $j \in [N]$, it samples the tuple $\mathsf{ptd}_j := (\mathbf U_j,\mathbf C_j,\mathbf{T}_{\mathbf{C}_j}) \leftarrow \mathsf{PTrapGen}(\mathbf C,\mathbf T_{\mathbf{C}},j)$. The algorithm outputs {\setlength{\abovedisplayskip}{2pt}
\setlength{\abovedisplayshortskip}{2pt}
\[
\mathsf{pk}:=(\mathsf{crs}_{\mathsf{samp}},\mathbf C,\mathbf u),
\qquad
\mathsf{sk}_j:=\mathsf{ptd}_j.
\]
}
\noindent$\mathsf{TBIBE.Enc}(\mathsf{pk}, r \in \mathcal{I}, \mathsf{m} \in \{0,1\}):$ The algorithm samples $s \leftarrow \mathbb{Z}_q^{nk}$ and set $\tilde{h} := H(r) \in \{0,1\}^{d}$. The matrix $\mathbf{A}_r := \mathsf{ExpandLocal} (1^\lambda,\mathsf{crs}_{\mathsf{samp}},\tilde{h}) \in \mathbb{Z}_q^{nk \times t}$ is then computed. It next selects $e_1 \leftarrow \mathcal{D}_{\mathbb{Z},\chi_1}, \mathbf e_2 \leftarrow \mathcal{D}_{\mathbb{Z},\chi_0}^{2mk} \text{ and } \mathbf e_3 \leftarrow \mathcal{D}_{\mathbb{Z},\chi_0'}^{t+\overline{m}}$. It then sets  $\mathbf F := [\,\mathbf A_r | H_{\mathsf ro}(r)\,] \in \mathbb{Z}_q^{nk \times (t + \overline{m})}$. The algorithm outputs the ciphertext

{\setlength{\abovedisplayskip}{2pt}
\setlength{\abovedisplayshortskip}{2pt}
\[
\mathsf{ct}:=(r,\mathsf{ct}[1],\mathsf{ct}[2],\mathsf{ct}[3])
\qquad \text{where} \qquad
\left\{
\begin{aligned}
\mathsf{ct}[1] &:= \mathbf{u}^{\top}\mathbf{s}+e_1+m\cdot \lfloor q/2 \rfloor \in \mathbb{Z}_q,\\
\mathsf{ct}[2] &:= \mathbf{C}^{\top}\mathbf{s}+\mathbf{e}_2 \in \mathbb{Z}_q^{2mk},\\
\mathsf{ct}[3] &:= \mathbf{F}^{\top}\mathbf{s}+\mathbf{e}_3 \in \mathbb{Z}_q^{t + \overline{m}}.
\end{aligned}
\right.
\]}
$\mathsf{TBIBE.PreDec}(\mathsf{sk_j}, T \subset_{t} [N], (r_i)_{i \in [\ell]}):$ For each $i \in [\ell]$, the pre-decryption algorithm sets $\tilde{h}_i := H(r_i) \in \{0,1\}^d$. The algorithm then check for the distinction among all the vector $\tilde{h}_1, \dots, \tilde{h}_\ell$, otherwise, it aborts. It then generates the matrices
{\setlength{\abovedisplayskip}{2pt}
\setlength{\abovedisplayshortskip}{2pt}\[
    (\mathbf A_1,\ldots,\mathbf A_\ell,\mathbf T_A)
    := \mathsf{Expand}(1^\lambda,1^\ell,\mathsf{crs}_{\mathsf{samp}},(\tilde{\mathbf h}_i)_{i\in[\ell]})  
    \]}

where \[
    \mathbf{A}
    :=
    \left[
    \begin{array}{cccc|c}
        \mathbf{A}_1 &              &        &              & \mathbf{G} \\
                     & \mathbf{A}_2 &        &              & \mathbf{G} \\
                     &              & \ddots &              & \vdots     \\
                     &              &        & \mathbf{A}_\ell & \mathbf{G}
    \end{array}
    \right]
    \in \mathbb{Z}_q^{\ell nk \times (\ell t+\tilde{m})}
\]

Let $
    \hat{\mathbf{A}}
    :=
    \begin{bmatrix}
        H_{\mathrm{ro}}(r_1) & & \\
        & H_{\mathrm{ro}}(r_2) & \\
        & & \ddots \\
        & & & H_{\mathrm{ro}}(r_\ell)
    \end{bmatrix}
    \in \mathbb{Z}_q^{\ell n k \times \ell \bar{m}} .
$ and employ $\mathbf{h} = (\mathbf{u},\dots,\mathbf{u})$, we calculate random coin 
{
$$\mathsf{rb}:=H_{\mathsf{sp}}(\mathbf A,\mathbf T_A,(r_i)_{i\in[\ell]},\mathbf u,\sigma_1)$$}

and invokes 
{\setlength{\abovedisplayskip}{2pt}
\setlength{\abovedisplayshortskip}{2pt} $$\operatorname{col}(\mathbf v_1^1,\ldots,\mathbf v_\ell^1,\hat{\mathbf c},\mathbf v_1^2,\ldots,\mathbf v_\ell^2)
    := \mathsf{SampleLeft}(\mathbf A,\hat{\mathbf A},\mathbf T_A,\mathbf h,\sigma_1,\mathsf{rb})
$$}
Finally, it derives the partial short pre-decryption key for each party $j \in [N]:$
{\setlength{\abovedisplayskip}{2pt}
\setlength{\abovedisplayshortskip}{2pt}
$$\mathbf{sbk}_j\leftarrow \mathsf{PSampPre}(\mathbf C,\mathsf{ptd}_j,T,\mathbf c:=\mathbf G\hat{\mathbf c},\sigma_2) \in \mathbb{Z}_{q}^{2mk}$$}

\noindent$\mathsf{TBIBE.Combine}(\mathsf{pk},T,(\mathsf{sbk}_i)_{i\in T}):$ The combination algorithm computes
    $$\mathbf{sbk} \leftarrow \mathsf{Rec}((\mathbf{sbk}_j)_{j \in T}) \in \mathbb{Z}_q^{2mk}$$ 
and return it as full-decryption key

\noindent$\mathsf{TBIBE.Dec}(\mathsf{pk},(\mathbf{sbk}), (\mathsf{ct_i})_{i \in [\ell]}):$ the decryption algorithm first also follow the same path as PreDec to caculate $\operatorname{col}(\mathbf v_1^1,\ldots,\mathbf v_\ell^1,\hat{\mathbf c},\mathbf v_1^2,\ldots,\mathbf v_\ell^2)$. Finally, for each $i \in [\ell]$, it set $\mathbf v_i:=\operatorname{col}(\mathbf v_i^1,\mathbf v_i^2)$ and output the message $\mathsf{m}'_i$: 
$$m_i' := \operatorname{round}\!\left(
\mathsf{ct}_i[1]-\mathsf{sbk}^{\top}\mathsf{ct}_i[2]-\mathbf v_i^{\top}\mathsf{ct}_i[3]
\right).$$
Here, $\mathsf{round}(x)$ outputs $1$ if $|x - \lfloor q/2 \rfloor| < \lfloor q/4 \rfloor$, and $0$ otherwise.

\begin{remark}
Both $\mathsf{TBIBE.PreDec}$ and $\mathsf{TBIBE.Dec}$ can work with smaller batches of size less than $\ell$. The algorithms remain unchanged in this case, except for working with smaller matrices.
\end{remark}
\section{Correctness}
\begin{theorem}[Correctness]
The TBIBE scheme is correct, as in Definition~\ref{def:TBIBE_correctness},
whenever $H$ is injective, as in Definition~\ref{def:hash}, and the
parameters are chosen as in Table~\ref{tab:parameters}.
\end{theorem}
\begin{proof}
Let $(r_i,m_i)_{i\in[\ell]}$ be a batch of identities and messages.
Since $H$ is injective, the labels
$
    H(r_1),\ldots,H(r_\ell)
$
are distinct. Hence, the shifted multi-preimage sampler is applied to
distinct labels and outputs vectors satisfying, for every $i\in[\ell]$,
\[
    \mathbf{F}_{r_i}\mathbf{v}_i+\mathbf{c}=\mathbf{u},
    \qquad
    \mathbf{F}_{r_i}:=
    [\mathbf{A}_i \mid H_{\mathsf{ro}}(r_i)].
\]
Moreover, by the correctness of the partial lattice trapdoor scheme,
the combined pre-decryption key
$
    \mathbf{sbk}
    \leftarrow
    \mathsf{Rec}\big((\mathbf{sbk}_j)_{j\in T}\big)
$
satisfies
\[
    \mathbf{C}\cdot\mathbf{sbk}=\mathbf{c}
    \pmod q
    \qquad\text{and}\qquad
    \|\mathbf{sbk}\|\le \beta .
\]
Therefore, for every $i\in[\ell]$,
$
    \mathbf{F}_{r_i}\mathbf{v}_i+\mathbf{C}\cdot\mathbf{sbk}
    =
    \mathbf{u}
    \pmod q .
$

Now consider the decryption expression. Using
\[
\begin{aligned}
    \mathsf{ct}_i[1]
    &=
    \mathbf{u}^{\top}\mathbf{s}_i
    +
    e_{1,i}
    +
    m_i\lfloor q/2\rfloor,\\
    \mathsf{ct}_i[2]
    &=
    \mathbf{C}^{\top}\mathbf{s}_i
    +
    \mathbf{e}_{2,i},\\
    \mathsf{ct}_i[3]
    &=
    \mathbf{F}_{r_i}^{\top}\mathbf{s}_i
    +
    \mathbf{e}_{3,i},
\end{aligned}
\]
we obtain
\[
\begin{aligned}
    \mathsf{ct}_i[1]
    &-
    \left(
        \mathbf{sbk}^{\top}\mathsf{ct}_i[2]
        +
        \mathbf{v}_i^{\top}\mathsf{ct}_i[3]
    \right)
    \\
    &=
    \mathbf{u}^{\top}\mathbf{s}_i
    +
    e_{1,i}
    +
    m_i\lfloor q/2\rfloor
    -
    \left(
        \mathbf{sbk}^{\top}\mathbf{C}^{\top}\mathbf{s}_i
        +
        \mathbf{sbk}^{\top}\mathbf{e}_{2,i}
        +
        \mathbf{v}_i^{\top}\mathbf{F}_{r_i}^{\top}\mathbf{s}_i
        +
        \mathbf{v}_i^{\top}\mathbf{e}_{3,i}
    \right)
    \\
    &=
    \left(
        \mathbf{u}
        -
        \mathbf{C}\mathbf{sbk}
        -
        \mathbf{F}_{r_i}\mathbf{v}_i
    \right)^{\top}
    \mathbf{s}_i
    +
    e_{1,i}
    -
    \mathbf{sbk}^{\top}\mathbf{e}_{2,i}
    -
    \mathbf{v}_i^{\top}\mathbf{e}_{3,i}
    +
    m_i\lfloor q/2\rfloor
    \\
    &=
    e_{1,i}
    -
    \mathbf{sbk}^{\top}\mathbf{e}_{2,i}
    -
    \mathbf{v}_i^{\top}\mathbf{e}_{3,i}
    +
    m_i\lfloor q/2\rfloor
    \pmod q .
\end{aligned}
\]

It remains to bound the total error term
$
    E_i
    :=
    e_{1,i}
    -
    \mathbf{sbk}^{\top}\mathbf{e}_{2,i}
    -
    \mathbf{v}_i^{\top}\mathbf{e}_{3,i}.
$
By item~\ref{param:w} and Lemma~\ref{lem:tail_bound},
$\mathbf{v}_i$ is distributed statistically close to a discrete Gaussian
preimage with parameter $\sigma_1$. Hence, except with negligible
probability,
\[
    \|\mathbf{v}_i\|
    \le
    \sigma_1\sqrt{t+\overline m}.
\]
Similarly, by the Gaussian tail bound and the parameter choices,
\[
    |e_{1,i}|\le \chi_1,
    \qquad
    \|\mathbf{e}_{2,i}\|\le \chi_0\sqrt{2mk},
    \qquad
    \|\mathbf{e}_{3,i}\|\le \chi_0'\sqrt{t+\overline m}.
\]
Using Cauchy--Schwarz, we get
\[
\begin{aligned}
    |E_i|
    &=
    \left|
        e_{1,i}
        -
        \mathbf{sbk}^{\top}\mathbf{e}_{2,i}
        -
        \mathbf{v}_i^{\top}\mathbf{e}_{3,i}
    \right|
    \\
    &\le
    |e_{1,i}|
    +
    \left|
        \mathbf{sbk}^{\top}\mathbf{e}_{2,i}
    \right|
    +
    \left|
        \mathbf{v}_i^{\top}\mathbf{e}_{3,i}
    \right|
    \\
    &\le
    |e_{1,i}|
    +
    \|\mathbf{sbk}\|\cdot\|\mathbf{e}_{2,i}\|
    +
    \|\mathbf{v}_i\|\cdot\|\mathbf{e}_{3,i}\|
    \\
    &\le
    \chi_1
    +
    \beta\chi_0\sqrt{2mk}
    +
    (t+\overline m)\sigma_1\chi_0'.
\end{aligned}
\]

By the modulus condition in Table~\ref{tab:parameters},
\[
    q
    >
    4\left(
        \chi_1
        +
        \beta\chi_0\sqrt{2mk}
        +
        (t+\overline m)\sigma_1\chi_0'
    \right).
\]
Hence,
$
    |E_i| < q/4
$
except with negligible probability. Hence,
\[
    \mathsf{ct}_i[1]
    -
    \left(
        \mathbf{sbk}^{\top}\mathsf{ct}_i[2]
        +
        \mathbf{v}_i^{\top}\mathsf{ct}_i[3]
    \right)
\]
is closer to \(0\) when \(m_i=0\), and closer to \(\lfloor q/2\rfloor\)
when \(m_i=1\). Thus the rounding procedure outputs
$
    m_i'=m_i
$
for every \(i\in[\ell]\), except with negligible probability. This proves
correctness.

Finally, the asymptotic sizes of the public key, pre-decryption key, and
ciphertext are
\[
\left\{
\begin{aligned}
|\mathsf{pk}|
&=
n(2t+m)\log q
=
O\!\left(
d\lambda^2\log(\lambda)\log(\ell)
+
d\lambda^2\log^2(\ell)
\right),
\\[0.4em]
|\mathsf{sbk}|
&=
O(k\sigma_2\sqrt{mk})
=
O\!\left(
\lambda\log^2\lambda
+
\lambda\log\lambda\log\ell
\right),
\\[0.4em]
|\mathsf{ct}|
&=
(1+2mk+t+\overline m)\log q.
\end{aligned}
\right.
\]
\end{proof}
% \end{document}